\section{Comparación controlada del aporte contextual}
\label{sec:met-controlado}

\subsection{Pregunta y tarea común}

El protocolo controlado pregunta si, en una tarea fija, conservar diferenciación entre ocurrencias mediante pseudo-sentidos o vectores originales modifica favorablemente los comportamientos globales y su robustez respecto de reemplazarla por el promedio por palabra. Se mantiene $K=80$, con las 640 palabras, ocho clases, etiquetas, competidores y 51 particiones del referente S. Cada palabra tiene una sola etiqueta, heredada por todas sus ocurrencias y pseudo-sentidos.

K80 se adopta porque el referente de ocho clases dispone de soporte estructural para los cinco bins en el recorrido auditado. La auditoría histórica del referente rechaza K40, K50, K60 y K70 por falta de cobertura, y registra K80 con 414 palabras FULL y 376 COMMON. La preparación deberá reproducir esos inventarios y su identidad antes de ejecutar los nuevos ensayos. No se comparará una tarea no estimable en un brazo con otra tarea distinta que sí lo sea, ni se elegirá $K$ por la magnitud de una ventaja observada.

Mantener ocho clases evita que dos copias del mismo promedio por palabra reciban etiquetas incompatibles procedentes de centros de pseudo-sentido distintos. Las categorías permanecen operacionales y comunes a todos los brazos. Este diseño responde a una pregunta condicionada a K80; no evalúa si inducir pseudo-sentidos permite construir tareas admisibles con un $K$ menor.

\subsection{Intervención A/P/O con posiciones comunes}

Para cada palabra $w$ y cada una de sus 100 posiciones canónicas se establecen tres representaciones:
\begin{equation}
\mathbf r_A(w,j)=\mathbf a_w,\qquad
\mathbf r_P(w,j)=\mathbf p_{w,g(j)},\qquad
\mathbf r_O(w,j)=\mathbf x_{wj},
\label{eq:met-intervencion}
\end{equation}
donde $g(j)$ es la pertenencia congelada de la ocurrencia a un pseudo-sentido. A elimina la diferenciación entre apariciones de la misma palabra; P conserva diferencias entre grupos, pero no dentro de cada grupo; O conserva los vectores individuales.

Todas las palabras aportan 100 posiciones en las tres representaciones. En A contienen copias de la media; en P, copias del prototipo según su soporte; en O, las ocurrencias originales. De esta manera, para cada split existen exactamente $100|\mathcal L_s|$ posiciones de memoria en cualquier brazo. Las consultas de inferencia conservan también las mismas posiciones e identificadores.

Las copias son una forma de controlar la multiplicidad del operador, no nuevas observaciones independientes. Eliminar duplicados sin compensación cambia los pesos de atención y puede cambiar de modo distinto Softmax, Entmax y Normmax. Una compresión computacional solo será válida si reproduce el operador con sus multiplicidades. También se verificará la identidad ponderada de la ecuación~\eqref{eq:met-identidad-medias}; no se construirá A como una media no ponderada de prototipos.

La intervención conserva simultáneamente el corpus, el encoder y las apariciones utilizadas. Las diferencias de norma, dispersión y geometría que resultan de promediar forman parte de la intervención representacional. Por ello, el efecto no se identifica exclusivamente con una noción lingüística abstracta de sentido.

\subsection{Nueve condiciones y contrastes}

Se cruzan tres representaciones de memoria con tres representaciones de consulta, según la tabla~\ref{tab:met-factorial}. Los contrastes principales son P/P frente a A/A y O/O frente a A/A. Los cruces permiten separar cambios de memoria y de consulta y examinar su interacción.

\begin{table}[tbp]
\centering
\caption{Diseño controlado de memoria por consulta. La primera letra identifica las memorias y la segunda las consultas de inferencia.}
\label{tab:met-factorial}
\begin{tabular}{@{}lccc@{}}
\toprule
Memoria / Consulta & A & P & O \\
\midrule
A & A/A & A/P & A/O \\
P & P/A & P/P & P/O \\
O & O/A & O/P & O/O \\
\bottomrule
\end{tabular}
\end{table}

Para cada tendencia se informarán, entre otras, las diferencias emparejadas
\begin{equation}
\Delta_k^{P}=T_k^{P/P}-T_k^{A/A},\qquad
\Delta_k^{O}=T_k^{O/O}-T_k^{A/A},\qquad k=1,2,3.
\label{eq:met-contrastes-control}
\end{equation}
Las comparaciones conservan palabra, ocurrencia, red, beta, split y semilla. Un cambio de signo esperado se distingue de un cambio de magnitud dentro del mismo signo. Si ambas condiciones cumplen una tendencia, ello no demuestra equivalencia; se comparan los márgenes y las omisiones críticas.

Los efectos de memoria pueden examinarse comparando P/A u O/A con A/A, y los de consulta comparando A/P o A/O con A/A. Para una tendencia $T_k$, una diferencia de diferencias como
\begin{equation}
\mathcal I_k^P=T_k^{P/P}-T_k^{P/A}-T_k^{A/P}+T_k^{A/A}
\label{eq:met-interaccion-control}
\end{equation}
resume la no aditividad de los dos cambios sobre esa escala. Se trata de un contraste descriptivo del diseño, no de una prueba inferencial que suponga independencia de las celdas o de los ensayos.

\subsection{Controles de aprendizaje, escala y clasificación}

Todas las ocurrencias y pseudo-sentidos de una palabra permanecen en aprendizaje o en inferencia según el split común. No se incorporan vectores de palabras de inferencia a las memorias, al entrenamiento del clasificador ni a la calibración. Sus promedios y agrupaciones solo construyen sus consultas dentro del benchmark de representaciones fijas. El procedimiento no pretende evaluar una extracción e inducción completas sobre un corpus externo no observado.

Para cada split se calculará
\begin{equation}
\rho_s=\operatorname{mediana}\{\|\mathbf a_w\|_2:w\in\mathcal L_s\},
\label{eq:met-rho-comun}
\end{equation}
con una entrada por palabra. El mismo valor normalizará las claves de A, P y O según la ecuación~\eqref{eq:met-keynorm}. Valores y consultas permanecerán RAW. La calibración HFY utilizará únicamente la memoria A de aprendizaje, con las multiplicidades reales, y una escala por red y split se mantendrá en todos los brazos. No se recalibrarán P y O para recuperar determinados signos o para igualar sus errores.

El observador de clasificación utilizará siempre una consulta de entrenamiento A por palabra: su promedio RAW, recuperado con la memoria A, P u O correspondiente. Esos estados, una fila por palabra y con las mismas etiquetas, entrenarán el ensemble. Para cada memoria, red, beta, split y semilla se entrenará un único ensemble que evaluará los tres tipos de consulta de inferencia.

Los coeficientes del clasificador podrán cambiar cuando cambien los estados recuperados. Lo que permanece común es la receta de entrenamiento, sus unidades, etiquetas y semillas; no se impone artificialmente el mismo clasificador ajustado a representaciones diferentes. La intervención mide el efecto de cambiar la representación bajo ese procedimiento común de observación.

\subsection{Evaluación común y conclusiones permitidas}

Los competidores, similitudes, momentos de clase, COMMON y cinco bins se conservarán a partir del referente S. No se reconstruirán vecinos para cada representación. La expansión a posiciones de ocurrencia no modifica el inventario geométrico de palabras. Las tasas se agregarán primero por palabra y luego con los mismos pesos de clase/bin en todos los brazos.

Se utilizarán CORE3, una iteración, 98 betas, 51 splits y diez semillas de clasificación. El modelo principal será \texttt{S\_TENSOR\_SHARED} con Ridge $10^{-4}$ y la sensibilidad será \texttt{H\_BETA\_LINEAR\_Z} con Ridge $0{,}1$. Se conservarán CV por palabra, MATCHED5, GRID30 y ocho omisiones empíricas y estadísticas. Las discrepancias entre superficies y modelos se informarán sin seleccionar una receta por sus signos.

Antes de ejecutar se verificará la identidad de las 100 ocurrencias, sus pertenencias, promedios, etiquetas, posiciones y roles; se reproducirá la geometría K80 y se comprobará la procedencia exclusivamente de aprendizaje de las escalas. Si un control falla, se documentará la incompatibilidad antes de modificar el diseño. No se sustituirán datos o parámetros silenciosamente.

La comparación permitirá atribuir diferencias a la intervención representacional definida dentro de esta tarea y este procedimiento. No demostrará superioridad para cualquier $K$, corpus o encoder, validez lingüística de todos los pseudo-sentidos ni causalidad psicológica humana. Los antecedentes exploratorios ya examinados forman parte del origen del diseño; la nueva ejecución no se presentará como una confirmación completamente independiente de toda la exploración previa. Los resultados nulos, contradictorios o desfavorables se conservarán.

